\documentclass[a4paper,11pt]{article}
\usepackage[utf8]{inputenc}
\usepackage[T1]{fontenc}
\usepackage[french]{babel}
\usepackage{amsmath,amsfonts,amssymb}
\usepackage{geometry}
\usepackage{fancyhdr}
\usepackage{titlesec}
\usepackage{enumitem}

\geometry{margin=2.5cm}

% Header and footer configuration
\pagestyle{fancy}
\fancyhf{}
\fancyhead[L]{Licence 1 Informatique}
\fancyhead[R]{CERI - Avignon Université}
\fancyfoot[C]{Mathématiques Discrètes \hfill TD5 - Prise en main de la combinatoire \hfill Page \thepage/3}

% Title formatting
\titleformat{\section}
{\Large\bfseries}
{\thesection}
{1em}
{}

\begin{document}

\begin{center}
{\Huge\bfseries TD5 - Prise en main de la combinatoire} \\[0.5cm]
{\large Hiver 2026} \\[0.3cm]
\end{center}

\vspace{0.5cm}

L'objectif de ce TD est de manipuler les concepts de base des ensembles et de la combinatoire.

\section{Application de la déduction logique aux ensembles}

\begin{enumerate}
\item Rédigez une argumentation utilisant les règles d'inférence pour montrer que si $A \subset B$ et $B \subset C$ alors on peut en déduire $A \subset C$ (où $A$, $B$ et $C$ sont des ensembles)

\item Rédigez une argumentation utilisant les règles d'inférence pour montrer que si $A \subset B$ et $B \subset C$ alors on peut en déduire $A \subset C$ (où $A$, $B$ et $C$ sont des ensembles)
\end{enumerate}

\section{Éléments d'ensembles}

Pour chaque ensemble ci-dessous, listez ses éléments.

\begin{enumerate}
\item $A = \{x \in \mathbb{R} \mid x^2 - 4 = 0\}$

\item $B = \{x \in \mathbb{Z} \mid x^2 = 2\}$

\item $C = \{x \in \mathbb{Z} \mid x^2 < 16\}$

\item $D = \{1, 2, 2, 3, 3, 3, 4, 4, 4, 4\}$

\item $E = \{x \in \mathbb{N} \mid x < 100 \wedge x = k^2, k \in \mathbb{N}\}$

\item $F = \{x \in \mathbb{N}^* \mid x < 10 \wedge x = 2 \cdot k + 1, k \in \mathbb{N}\}$
\end{enumerate}

\section{Manipulation de base des ensembles}

On considère la liste d'ensembles suivante :

\begin{minipage}[t]{0.45\textwidth}
\begin{align*}
\text{A} &= \{3, -3\} \\
\text{B} &= \{x \in \mathbb{Z} \mid x^2 = 5\} \\
\text{C} &= \{\{\emptyset\}\} \\
\text{D} &= \{-3, -3, 3, 3, 3, 3\} \\
\text{E} &= \{x \in \mathbb{Z} \mid x \geq 0\} \\
\text{F} &= \{x \in \mathbb{Z} \mid x^2 = 9\}
\end{align*}
\end{minipage}
\hfill
\begin{minipage}[t]{0.45\textwidth}
\begin{align*}
\text{G} &= \{x \in \mathbb{Z} \mid x^2 < 0\} \\
\text{H} &= \mathbb{N} \\
\text{I} &= \{x \in \mathbb{Z} \mid x^2 \leq 0\} \\
\text{J} &= \{-3, 3, \{3\}, \{-3\}\}
\end{align*}
\vspace{0.5cm}
\begin{align*}
\text{K} &= \{x \in \mathbb{N} \mid x^2 = 9\} \\
\text{L} &= \{x \in \mathbb{R} \mid x^3 = 9\} \\
\text{M} &= \{x \in \mathbb{Z} \mid x^3 + 3 = 3\} \\
\text{N} &= \emptyset \\
\text{O} &= \{3, -3, 3, 3, -3\}
\end{align*}
\end{minipage}

\begin{enumerate}
\item Lesquels de ces ensembles sont égaux ?

\item Lesquels de ces ensembles sont des sous-ensembles d'un autre ensemble de la liste ?

\item Quelle est la cardinalité de ces ensembles ?
\end{enumerate}

\section{Opérations dans les ensembles}

Pour chaque paire d'ensembles $A$ et $B$ suivantes, donnez les ensembles correspondant à $A \cup B$, $A \cap B$, $A \setminus B$, $B \setminus A$, ainsi que $\overline{A}$ et $\overline{B}$ dans le cas où $A$ et $B$ sont donnés.

\begin{enumerate}
\item $A = $ ``Ensemble des étudiants habitant à moins de 3km du CERI.''\\
$B = $ ``Ensemble des étudiants qui viennent au CERI à pieds.''\\
Où $A$ et $B$ sont dans l'ensemble des étudiants du CERI

\item $A = \{1, 2, 3, 4, 5\}$\\
$B = \{0, 3, 6\}$\\
Où $A$ et $B$ sont dans $\{0, 1, 2, 3, 4, 5, 6, 7, 8, 9, 10\}$

\item $A = \{a, b, c, d, e\}$\\
$B = \{a, b, c, d, e, f, g, h\}$\\
Où $A$ et $B$ sont dans l'ensemble $\{a, b, c, d, e, f, g, h\}$

\item $A = \{x \in \mathbb{N} \mid x = 2 \cdot k, k \in \mathbb{N}\}$\\
$B = \{x \in \mathbb{N} \mid x < 100\}$\\
Où $A$ et $B$ sont dans l'ensemble $\mathbb{N}$

\item $A = $ ``Ensemble des étudiants de licence informatique''\\
$B = $ ``Ensemble des étudiants qui suivent des cours d'algèbre ou des cours d'analyse.''\\
Où $A$ et $B$ sont dans l'ensemble des étudiants d'Avignon Université
\end{enumerate}

\section{Application des principes de bases}

\begin{enumerate}
\item Un groupe d'étudiants contient 325 étudiants d'informatique et 18 étudiants de mathématiques.
\begin{enumerate}
\item Combien y a-t-il de manières de choisir un représentant qui soit étudiant d'informatique et l'autre soit étudiant de mathématiques ?
\item Combien y a-t-il de manières de choisir deux représentants qui soit étudiant d'informatique ou étudiant de mathématiques ?
\end{enumerate}

\item Un QCM contient 4 questions. Chaque question propose 4 réponses possibles.
\begin{enumerate}
\item Combien de manière y a-t-il de répondre QCM si on donne une réponse par question ?
\item Combien de manière y a-t-il de répondre QCM si on peut donner une réponse ou ne pas répondre ?
\end{enumerate}

\item Un groupe d'étudiants contient des étudiants de licence informatique, licence mathématiques et des étudiants en double-licence. Il y a 87 étudiants en informatique (licence seule ou double-licence), 23 étudiants qui font une licence de mathématiques (licence seule ou double-licence) et 7 étudiants qui font une double-licence. Combien d'étudiants y a-t-il dans le groupe ?
\end{enumerate}

\section{Photos de mariage}

Alix et Camille se marient et veulent faire des photos avec leurs 4 plus proches invités. La photographe propose d'aligner ces personnes sur les photos. Combien y a-t-il de possibilité d'arranger ces 6 personnes selon les contraintes suivantes :

\begin{enumerate}
\item Aucune contrainte

\item Alix et Camille doivent être côte à côte

\item Alix et Camille ne doivent pas être côte à côte

\item Alix doit se situer quelque part sur la gauche de Camille (mais pas forcément à côté)

\item Alix doit se situer quelque part sur la droite de Camille (mais pas forcément à côté)
\end{enumerate}

\section{Boutique de vêtements}

Une boutique de vêtement lance une nouvelle collection de t-shirts. Un t-shirt est défini par sa coupe, sa taille, sa couleur, son motif. On donne les ensembles de motifs, couleurs, coupes et taille :

\begin{itemize}
\item L'ensemble des coupes possibles est : coupe homme, coupe femme, coupe neutre
\item L'ensemble des tailles est : XS, S, M, L, XL, XXL
\item L'ensemble des couleurs possibles est : noir, bleu, rouge, vert, jaune, violet, marron
\item L'ensemble des motifs possibles est : uni, pois, rayures, dessin
\end{itemize}

\begin{enumerate}
\item Définissez l'ensemble des t-shirts. Quelle est la taille de cet ensemble ?

\item Combien de possibilités de t-shirts a-t-on dans les conditions ? Utilisez un diagramme en arbre pour répondre.
\end{enumerate}

Supposons à présent que toutes les combinaisons ne sont pas possibles. Les tailles XL et XXL ne sont disponibles que pour la coupe homme. Les couleurs noir, bleu et rouge sont disponibles pour toutes les coupes. Les tailles XS et S ne sont qu'en coupe femme. Les couleurs vert et jaune ne sont disponibles que pour les tailles M et L. Les couleurs violet et marron ne sont disponibles que pour les tailles XXL. Les couleurs violet et marron ne sont disponibles que pour les tailles M et L. Le motif à dessin est disponible uniquement pour les tailles XS, S et M. Le motif uni est disponible uniquement pour les tailles XS, S, M et L. Le motif pois est disponible uniquement pour les tailles M, XL et XXL. Le motif à rayures est disponible uniquement pour les tailles XS et S.

\section{Chaînes de caractères}

On s'intéresse ici aux chaînes de 8 caractères, où les seuls caractères considérés sont les 26 lettres de l'alphabet en minuscule. On rappelle que les lettres $a$, $e$, $i$, $o$, $u$ et $y$, et les consonnes sont les autres lettres.

\begin{enumerate}
\item Combien de chaînes de caractères sont possibles si aucun critère n'est présent et qu'on peut répéter des lettres ?

\item Combien de chaînes de caractères sont possibles si aucun voyelle n'est présente et qu'on ne peut pas répéter de lettres ?

\item Combien de chaînes de caractères commençant par une voyelle sont possibles si on peut répéter des lettres ?

\item Combien de chaînes de caractères commençant par une voyelle sont possibles si on ne peut pas répéter des lettres ?

\item Combien de chaînes de caractères contenant au moins une voyelle sont possibles si on peut répéter des lettres ?

\item Combien de chaînes de caractères contenant exactement une voyelle sont possibles si on peut répéter des lettres ?

\item Combien de chaînes de caractères commençant par un $x$ et contenant au moins une voyelle sont possibles si on peut répéter des lettres ?

\item Combien de chaînes de caractères commençant et finissant par un $x$ et contenant au moins une voyelle sont possibles si on peut répéter des lettres ?
\end{enumerate}

\end{document}
